\section{Conclusions}
\label{sec:conclusions}

\begin{enumerate}[leftmargin=1.6em,itemsep=4pt,topsep=4pt,parsep=0pt]

\item \emph{What is new.} In the millimetre band the record for nearby stars
was two measurements of one spectral line, each no more than
\SbPriorBandMHz\,MHz around \SbPriorFreqGHz\,GHz with no drift search and no
test of persistence \citep{Steffes1994,Mauersberger1996}. The carrier search
here is \NWinA{} fine-channel windows toward \NSysClassA{} systems within
40\,pc, \CvUnionA\,GHz of unique frequency in \CovNIslands{} disjoint intervals
and at least \SbPriorRatio{} times that bandwidth, of which \EdrNWinA{} windows
carry a track long enough for the searched drift to cross a channel; a
coarse-channel search for a steady unresolved excess adds \NWinB{} windows
toward \PrSysClassB{} systems. Both come from one pass through \NEB{} public
ALMA execution blocks spanning \SurvFreqLo--\SurvFreqHi\,GHz
(\S\ref{sec:sample}).

\item \emph{How deep.} Taken over the bandwidth the line mask leaves, the
median system's best window would have produced a threshold crossing from a
carrier of $\EirpNinetySysMedA$\,W nine times in ten, and \NSysEfifteen{} of
the \NSysALim{} systems with a constrained window reach $10^{15}$\,W;
combining each star's epochs in its own rest frame deepens the flux-density
limit by a median factor \SyStkGainMed, the deepest being Proxima Centauri's
$\StkProxEirpCal$\,W in coarse \ChanBModeMHz\,MHz channels, which bounds a
steady unresolved excess and not a drifting carrier (\S\ref{sec:stack}). That
is the completeness of a single observation and not of a confirmation, and for
\SmcNEstSys{} of the \NSysALim{} systems it is carried across from the
windows injected into rather than measured in their own. It also leaves out
atmospheric decorrelation, which a real carrier suffers and an injected one
does not, so the net bias on any one limit runs either way
(\S\ref{sec:pninetysel}).

\item \emph{What was found.} The search produced \EvNCrossRaw{} threshold
crossings -- \CkNCrossCenA{} Class~A in the primary census,
\CeNTailCrossWord{} Class~A in archival-tail blocks that enter no expectation,
and \PrFlagXrossWord{} Class~B, all of them from the single execution block
that fails the data-quality criterion of Appendix~\ref{sec:blockquality} -- so
the quality-passing data hold \EvNCrossQp{} crossings, every one from the
carrier search. \EvNAttrQp{} of those are attributed to a catalogued molecular
transition within $\pm\MaskHalfKms$\,km\,s$^{-1}$ of the star's velocity,
circumstellar or interstellar, and \EvNUnattrQp{} are not; an attribution is
not an identification, and \S\ref{sec:mmdifferent} says which carry evidence
beyond the coincidence. Of the unattributed, \CeTlObsCutWord{} is atmospheric,
leaving \CeObs{} of the primary census against the \CeExp{} these fields' own
control positions predict at probability \CeProbNull{} -- consistency, and not
a demonstration that the crossings are noise, since a further \CeAddLim{} would
have been consistent too (\S\ref{sec:trigunif}).

\item \emph{What the repeat observations show.} No crossing satisfies the
confirmation criterion. A transmitter on an orbiting body arrives at a
frequency its own motion has moved, so the test that bears on one is the broad
one, and every unattributed crossing another block covers was searched twice
over that way: across the whole covering window at every drift of the grid,
where \WuNWinExc{} of the \WuNWin{} distinct covering windows hold a statistic
above the trigger against the \WuExpWin{} chance predicts ($P=\WuProbWin$),
each of them a window whose crossing the ledger already lists; and pair by
pair, where of the \TjNPair{} pairs of crossings carrying enough measurement to
be asked the searched grid separates \TjNExcl{} and leaves \TjNOpen{}
unresolved, \TjNOpenPlain{} \emph{open} and \TjNPairEdge{}
\emph{edge-limited}. Neither establishes a recurrence and neither excludes one.
Nothing returned to the single cell a crossing's own discovery predicts either,
the largest single-block reading at the carried drift being \RcTMatched{} and
\RcTMax{} the largest at any drift of the grid, against a trigger of
\StageNullTrig; the same reading recovers \MkAttrRecur{} of the \MkAttrTested{}
line-attributed crossings it was tested on, each in a different observing
session more than a day away, so the machinery does return a repeat where one
exists. What that reading tests is narrow, and the narrowness belongs to the
criterion and not to the archive: it requires a carrier radiating in both
observations whose stellar-frame frequency comes back to within half a native
channel, and a transmitter accelerating at the \ZtALosLo{} to at least
\ZtALosHi\,m\,s$^{-2}$ every one of the \RcNExcl{} shows re-enters that
interval only if its motion is periodic and the second epoch catches it at the
matching phase, so this is not a test of a persistent transmitter of any kind.
Of the \NSysClassA{} systems covered, \ZcConfSysDay{} were observed again in
another observing session more than a day later at a frequency the first block
covered and \ZcConfSysYr{} more than a year later, while for the remaining
\ZcNoRepSys{} no repeat exists, so no transmitted power would have made a
confirmation possible: coverage is one limit on confirmation and the criterion
applied to the repeat is the other (\S\ref{sec:mmdifferent},
\S\ref{sec:etacrv}).

\item \emph{What this does not constrain.} These limits hold only inside the
domain actually searched: the drift grid stops at \DriftCeilLo\,nHz,
\AccelSheikhPct{} per cent of the ceiling recommended for a general carrier
search \citep{Sheikh2019}, so an unbound transmitter, one on a rapidly rotating
planet, or one whose orbital and rotational accelerations together exceed that
is not constrained here at all (\S\ref{sec:excludes}). The obstacle transfers
beyond this archive: channelisation chosen by other observers, and not
collecting area, decides which search is possible. Three choices a proposal can
simply make would lift it -- fine channels everywhere, frequencies and targets
chosen against the line forest rather than inherited from programmes that
sought it, and repeat visits at one tuning days and years apart
(\S\ref{sec:future}).

\end{enumerate}

%% ★★★★ ROUND 25: EIGHT NUMBERED CONCLUSIONS BECAME FIVE, to R1's "no more
%% than five short paragraphs".  Merged: old 1+2 (what is new / what was
%% searched) into 1; old 7 (stacking) into 2, keeping the median gain and
%% Proxima's deepest figure and dropping the inclination arithmetic, which
%% Sec. 5.7 carries; old 5+6 (repeat exclusions / confirmability) into 4.
%% Item 4 runs to six sentences rather than three: it carries two merged
%% items and A1's reframing, and R1's binding constraint is the FIVE
%% paragraphs, not a sentence count.
%%
%% ★★★★ THIS FILE FOLDS IN TWO OTHER AGENTS' HUNKS, because a whole-file
%% record would otherwise clobber them and their own `old` strings are
%% anchored in the PRISTINE eight-item structure that no longer exists:
%%  * r25-recur (recur.patch, both hunks on this file): A1's reframing.  The
%%    broad window search and the PAIR TEST become the headline recurrence
%%    result and the single-cell reading is demoted to "Nothing returned ...
%%    either"; the item is retitled "What the repeat observations SHOW" and
%%    not "exclude"; \TjNOpenPlain and \TjNPairEdge are theirs; and old item
%%    6's close becomes "coverage is one limit on confirmation and the
%%    criterion applied to the repeat is the other".
%%  * r25-cut-app (patches_cutapp.json X3.1): "(Table~\ref{tab:confirm})" is
%%    removed, Table D3 having moved to the deposit.
%% ►► IF r25-recur REVISES ITS TEXT AFTER THIS RECORD WAS WRITTEN, THIS FOLD
%%    IS STALE AND THEIRS WINS -- re-read before integrating.
%%
%% ★★★ RULING 11.2 RE-POINTS censusacct_r17's A4 AT SEC. 5.3, AND I HAVE
%% DELIBERATELY NOT TAKEN THE LICENCE IT OFFERS.  A4 required THIS file to
%% carry \CkNCrossCenA, \PrFlagXrossWord, "\CeNTailCrossWord{} Class~A",
%% \EvNCrossRaw and \EvNCrossQp, and the sentence carrying \PrFlagXrossWord
%% also to say "fails the data-quality criterion" and name a "single execution
%% block".  Item 3 above still satisfies A4 in ONE sentence instead of four,
%% for two reasons that are measurements and not caution:
%%  (1) \PrFlagXrossWord is cited in exactly TWO files, this one and
%%      sections/abstract.tex.  r25-abs is rewriting the abstract to 180-200
%%      words; if that drops it and this file drops it too, retire_macros
%%      deletes the macro and the flagged-block count leaves the paper
%%      silently.  Keeping it here costs one clause and removes that coupling.
%%  (2) The record therefore applies whether or not the A4 re-point has landed
%%      in the integration tree, which a record that is the only surviving
%%      copy of the work ought to do.
%% ► A fuller cut IS available and is the integrator's to take: delete the
%%   dash-clause naming the three sets from item 3 once (a) the A4 re-point is
%%   in the tree and (b) \PrFlagXrossWord is confirmed still cited elsewhere.
%%
%% ★★ WHAT WENT, AND WHERE IT LIVES.  The narrower-field reading
%% (\FrNarrowObs, \FrNarrowExp, \FrNarrowProb, \FrNarrowAddLim) is no longer
%% quoted here: the coordinator's A2 ruling puts the alternative reading in
%% Sec. 5.3.1 and Appendix D as a labelled sensitivity test, and B9 asks for it
%% to be quoted ONCE.  The primary reading kept here is the one that is NOT
%% significant, so dropping the alternative cannot strengthen a claim.
%% Proxima's inclination condition (\PxNChanWalk, \PxIncMaxDeg,
%% \PxEirpSingle) stays in Sec. 5.7, where the stack's one-channel tolerance is
%% derived; the conclusion above no longer quotes the stacked figure as
%% applying to the planet, so the qualification it needed is not missing.
